\section{Introduction}
\label{sec:intro}

Docking gives a resident underwater vehicle a place to recharge and offload data without surfacing, and it removes the slow, weather-limited and expensive recovery by a support vessel from the mission cycle. The risk concentrates in the final metres, where the vehicle flies onto the station on its forward camera. Real stations are not still. A tether management system suspended from a vessel follows the vessel's heave: in the only reported autonomous docking to such a cage, the cage moved 1.1~m peak to peak at 8.5~s \cite{trslic2020vision}. A moored station follows the wave orbital motion at its depth, which decays exponentially with depth at a rate set by the wave period, so long swell reaches a dock at continental-shelf depth with displacements of order 0.1~m while short wind sea does not \cite{holthuijsen2007waves}. Off New South Wales, the operating context of this work, the mean sea state is 1.6~m at 8~s \cite{short1992wave,kinsela2024nearshore}, which forces a 0.1~m sway on a dock near 33~m depth.

The docking control literature treats the dock as a fixed pose and its estimator as a noise filter. Trslic et al.\ servoed reactively on the latest measured pose and accepted contact during dynamic docking \cite{trslic2020vision}; their later heave predictor used a pressure sensor on the cage and was evaluated offline \cite{trslic2020neuro}. Recent marker-based systems fuse the camera pose with an extended Kalman filter but keep the dock frame as the world frame \cite{gunzel2026docking,ni2025vision}. The obvious extension for a moving station is to estimate the dock velocity onboard and feed it forward into the vehicle command, so that the vehicle rides along with the dock and the position loop only trims the residual. To our knowledge no prior work closes that loop through vision alone.

This paper reports what happens when the obvious extension is added to a passive-marker terminal docking stage for a low-cost ROV. In simulation with the production autopilot in the loop, a reactive baseline docks cleanly onto a station swaying with 0.1~m amplitude down to 8~s periods, whereas a constant-velocity estimator with velocity feedforward degrades at 10~s and fails at 8 and 6~s. The estimator's own signals point at itself: its velocity state reaches many times the true dock amplitude in the failed trials. Rebuilding every measurement from the recorded trials shows that this reading is wrong, and the paper's contributions follow from what the data show instead:
\begin{enumerate}
  \item a failure mechanism established from the recorded trials rather than from the system's own diagnostics: the dock measurement and the velocity estimate are accurate, the feedforward over-drives a plant whose gain and lag differ from those it was tuned to, and the estimator's apparent runaway is dead-reckoning during the marker blackout that the over-swing causes;
  \item a redesigned feedforward that closes on the vehicle's own velocity, with a bounded velocity state during blackout, evaluated against the reactive baseline, the original method and an oracle feedforward from ground-truth dock velocity, which separates estimator-limited from plant-limited failure;
  \item an observability result for the estimator under correlated navigation error: a camera measuring relative position cannot distinguish a slowly varying vehicle-velocity error from dock velocity, for a world-frame or a relative-frame estimator alike, and only the fast part of the error is removed by estimating in the relative frame.
\end{enumerate}

Section~\ref{sec:related} places the work, Section~\ref{sec:system} describes the system, Section~\ref{sec:failure} establishes the mechanism, Section~\ref{sec:redesign} gives the redesign, Section~\ref{sec:design} the experiments, and Sections~\ref{sec:results} to~\ref{sec:conclusion} the results and what they mean for a deployment.
